\chapter{Unit 2: Fundamentals of A.C. Circuits}

\section{Sinusoidal Signal Measures and Definitions}

\begin{formulabox}{Mathematical Representation of Alternating Signals}
\[ v(t) = V_m \sin(\omega t + \theta), \qquad i(t) = I_m \sin(\omega t + \phi) \]
\begin{itemize}[leftmargin=1.5em]
    \item \textbf{Peak Amplitude}: $V_m$ (maximum excursion from zero).
    \item \textbf{Peak-to-Peak Amplitude}: $V_{p-p} = 2 V_m$.
    \item \textbf{Angular Frequency}: $\omega = 2\pi f = \frac{2\pi}{T}$ [rad/s].
    \item \textbf{Cyclic Frequency}: $f = \frac{1}{T}$ [Hz], where $T$ is the period in seconds.
    \item \textbf{Phase Difference}: $\Delta \phi = \theta - \phi$. If $\Delta \phi > 0$, voltage \textbf{leads} current by $\Delta \phi$. If $\Delta \phi < 0$, voltage \textbf{lags} current.
\end{itemize}
\end{formulabox}

\begin{formulabox}{RMS, Average, Form Factor, and Peak Factor}
\begin{itemize}[leftmargin=1.5em]
    \item \textbf{Root Mean Square (RMS / Effective) Value}: The DC equivalent heating value:
    \[ V_{rms} = \sqrt{\frac{1}{T}\int_0^T v^2(t) dt} \implies \text{For pure sine wave: } V_{rms} = \frac{V_m}{\sqrt{2}} \approx 0.707 V_m \]
    \item \textbf{Average (Mean) Value} (taken over half-cycle for symmetric AC):
    \[ V_{avg} = \frac{1}{T/2}\int_0^{T/2} v(t) dt \implies \text{For pure sine wave: } V_{avg} = \frac{2 V_m}{\pi} \approx 0.637 V_m \]
    \item \textbf{Form Factor ($k_f$)}: Ratio of RMS to Average value:
    \[ k_f = \frac{V_{rms}}{V_{avg}} = \frac{V_m / \sqrt{2}}{2 V_m / \pi} = \frac{\pi}{2\sqrt{2}} \approx 1.11 \quad \text{(Pure Sine Wave)} \]
    \item \textbf{Peak / Crest Factor ($k_p$)}: Ratio of Peak to RMS value:
    \[ k_p = \frac{V_m}{V_{rms}} = \frac{V_m}{V_m / \sqrt{2}} = \sqrt{2} \approx 1.414 \quad \text{(Pure Sine Wave)} \]
\end{itemize}
\end{formulabox}

\section{Phasors and Complex Impedance}

\begin{formulabox}{Complex Impedance and Admittance}
\begin{itemize}[leftmargin=1.5em]
    \item \textbf{Phasor Representation}: $v(t) = V_m \cos(\omega t + \theta) \iff \mathbf{V} = V_{rms}\angle \theta = V_{rms}e^{j\theta}$.
    \item \textbf{Ohm's Law in Frequency Domain}: $\mathbf{V} = \mathbf{I}\mathbf{Z} \iff \mathbf{I} = \mathbf{V}\mathbf{Y}$.
    \item \textbf{Impedance}: $\mathbf{Z} = R + jX = |\mathbf{Z}|\angle \theta_z$ [Ohms, $\Omega$]:
    \begin{align*}
        \text{Resistor: } & \mathbf{Z}_R = R \angle 0^\circ = R + j0 \\
        \text{Inductor: } & \mathbf{Z}_L = j\omega L = \omega L \angle 90^\circ \implies X_L = \omega L = 2\pi f L \\
        \text{Capacitor: } & \mathbf{Z}_C = \frac{1}{j\omega C} = -j\frac{1}{\omega C} = \frac{1}{\omega C} \angle -90^\circ \implies X_C = \frac{1}{\omega C} = \frac{1}{2\pi f C}
    \end{align*}
    \item \textbf{Impedance Magnitude and Phase}:
    \[ |\mathbf{Z}| = \sqrt{R^2 + (X_L - X_C)^2}, \qquad \theta_z = \arctan\left(\frac{X_L - X_C}{R}\right) \]
    \item \textbf{Admittance}: $\mathbf{Y} = \frac{1}{\mathbf{Z}} = G + jB$ [Siemens, S] ($G$: Conductance, $B$: Susceptance).
\end{itemize}
\end{formulabox}

\section{Series Resonance in RLC Circuits}

\begin{formulabox}{Series Resonance Master Relations}
Resonance occurs when the net reactance vanishes ($X_L = X_C$), meaning impedance is purely resistive ($\mathbf{Z} = R$) and current is at its maximum:
\begin{itemize}[leftmargin=1.5em]
    \item \textbf{Condition for Resonance}: $\omega_0 L = \frac{1}{\omega_0 C}$.
    \item \textbf{Resonant Frequency}:
    \[ \omega_0 = \frac{1}{\sqrt{LC}} \text{ [rad/s]}, \qquad f_0 = \frac{1}{2\pi\sqrt{LC}} \text{ [Hz]} \]
    \item \textbf{Impedance at Resonance}: $\mathbf{Z}_0 = R$ (minimum value). \quad \textbf{Current}: $I_0 = \frac{V}{R}$ (maximum value).
    \item \textbf{Power Factor at Resonance}: $\cos\phi = 1.0$ (\textbf{Unity}).
    \item \textbf{Quality Factor ($Q_0$)} (Magnification factor):
    \[ Q_0 = \frac{\omega_0 L}{R} = \frac{1}{\omega_0 C R} = \frac{1}{R}\sqrt{\frac{L}{C}} = \frac{V_L}{V} = \frac{V_C}{V} \]
    \item \textbf{Bandwidth ($BW$)} between half-power ($-3\text{ dB}$) points:
    \[ BW = \Delta f = f_2 - f_1 = \frac{f_0}{Q_0} = \frac{R}{2\pi L} \text{ [Hz]}, \qquad \Delta \omega = \frac{\omega_0}{Q_0} = \frac{R}{L} \text{ [rad/s]} \]
    \item \textbf{Half-Power Frequencies}:
    \[ \omega_1 = \omega_0 - \frac{R}{2L} + \sqrt{\omega_0^2 + \left(\frac{R}{2L}\right)^2}, \qquad \omega_2 = \omega_0 + \frac{R}{2L} + \sqrt{\omega_0^2 + \left(\frac{R}{2L}\right)^2} \]
    For high $Q$ ($Q \ge 10$): $f_1 \approx f_0 - \frac{BW}{2}, \quad f_2 \approx f_0 + \frac{BW}{2}, \quad f_0 = \sqrt{f_1 f_2}$.
\end{itemize}
\end{formulabox}

\section{Power Calculations and Power Factor}

\begin{formulabox}{Power Triangle: Real, Reactive, Apparent, and Complex Power}
\begin{itemize}[leftmargin=1.5em]
    \item \textbf{Real / Active Power ($P$)}: $P = V I \cos\phi = I^2 R$ [Watts, W].
    \item \textbf{Reactive Power ($Q$)}: $Q = V I \sin\phi = I^2 X$ [Volt-Amperes Reactive, VAR].
    \begin{itemize}
        \item Inductive load: $Q > 0$ (absorbs lagging VARs).
        \item Capacitive load: $Q < 0$ (delivers leading VARs).
    \end{itemize}
    \item \textbf{Apparent Power ($S$)}: $S = V I = I^2 |\mathbf{Z}|$ [Volt-Amperes, VA].
    \item \textbf{Complex Power ($\mathbf{S}$)}:
    \[ \mathbf{S} = \mathbf{V}\mathbf{I}^* = P + jQ = S \angle \phi \]
    \item \textbf{Power Triangle Relation}: $S = \sqrt{P^2 + Q^2}$.
    \item \textbf{Power Factor ($pf$)}: $\cos\phi = \frac{P}{S} = \frac{R}{|\mathbf{Z}|}$.
    \item \textbf{Power Factor Correction Capacitor}:
    To improve power factor from $\cos\phi_1$ to $\cos\phi_2$ for active load $P$:
    \[ Q_C = P(\tan\phi_1 - \tan\phi_2) \implies C = \frac{Q_C}{\omega V^2} = \frac{P(\tan\phi_1 - \tan\phi_2)}{\omega V^2} \]
\end{itemize}
\end{formulabox}

\section{Three-Phase Balanced Circuits}

\begin{formulabox}{Three-Phase Star and Delta Connections}
\begin{itemize}[leftmargin=1.5em]
    \item \textbf{Star (Wye, Y) Connection}:
    \[ V_L = \sqrt{3} V_{ph} \angle +30^\circ, \qquad I_L = I_{ph} \]
    Neutral current for balanced load: $I_N = I_A + I_B + I_C = 0$.
    \item \textbf{Delta ($\Delta$) Connection}:
    \[ V_L = V_{ph}, \qquad I_L = \sqrt{3} I_{ph} \angle -30^\circ \]
    \item \textbf{Total Balanced Three-Phase Power}:
    \begin{align*}
        P_{3\phi} &= \sqrt{3} V_L I_L \cos\phi = 3 V_{ph} I_{ph} \cos\phi \quad \text{[Watts]} \\
        Q_{3\phi} &= \sqrt{3} V_L I_L \sin\phi = 3 V_{ph} I_{ph} \sin\phi \quad \text{[VAR]} \\
        S_{3\phi} &= \sqrt{3} V_L I_L = 3 V_{ph} I_{ph} \quad \text{[VA]}
    \end{align*}
    \item \textbf{Two-Wattmeter Method}:
    \begin{align*}
        W_1 &= V_L I_L \cos(30^\circ - \phi), \qquad W_2 = V_L I_L \cos(30^\circ + \phi) \\
        P_{total} &= W_1 + W_2 \\
        \tan\phi &= \sqrt{3}\frac{W_1 - W_2}{W_1 + W_2} \implies \cos\phi = \frac{1}{\sqrt{1 + 3\left(\frac{W_1 - W_2}{W_1 + W_2}\right)^2}}
    \end{align*}
    \begin{itemize}
        \item If $W_1 = W_2 \implies \phi = 0^\circ \implies \cos\phi = 1.0$ (Unity pf).
        \item If one wattmeter reads zero ($W_2 = 0$) $\implies \phi = 60^\circ \implies \cos\phi = 0.5$.
        \item If one wattmeter reads negative $\implies \phi > 60^\circ \implies \cos\phi < 0.5$.
    \end{itemize}
\end{itemize}
\end{formulabox}
